\documentclass[11pt]{article}

\usepackage[margin=1in]{geometry}
\usepackage{amsmath}
\usepackage{amssymb}
\usepackage{booktabs}
\usepackage[T1]{fontenc}
\usepackage[utf8]{inputenc}
\usepackage{lmodern}

\title{Reactor Kinetics Lab Physics and Numerical Methods}
\author{}
\date{}

\begin{document}

\maketitle

\section{Purpose}

This note is the project-wide physics and numerics reference for the current
\emph{Reactor Kinetics Lab} repository.  It explains how the implemented models
are split across:
\begin{itemize}
  \item the OpenMC MGXS export notebook \texttt{openmc/exportMGXS.ipynb},
  \item the resolved multigroup diffusion notebook
        \texttt{theory/concentricModel.ipynb},
  \item the exploratory one-group notebook \texttt{theory/reactorModel.ipynb},
\end{itemize}

The notebooks are important \textbf{calibration and derivation workspaces}, while
the Python backend and React frontend are the \textbf{runtime implementation}.
The most important architectural points are:

\begin{quote}
The repository now contains \textbf{two coupled modeling tracks}: an
\textbf{OpenMC-informed resolved multigroup diffusion} pipeline that generates
transport-derived, cell-homogenized group constants from Monte Carlo and solves
the diffusion eigenvalue problem while preserving the exported cell regions,
and an \textbf{interactive web application} with a fast
\textbf{point-kinetics surrogate} on the Overview page and the validated
OpenMC-informed multigroup result exposed as the Core page.  The older
one-group Core and Transient Diffusion visualizations remain backend reference
models only; they are no longer frontend pages.
\end{quote}

\section{Role of the notebook}

The notebook \texttt{theory/reactorModel.ipynb} studies a heavy-water moderated
annular reactor with cylindrical symmetry in $(r,z)$ and a top-inserted central
control/shutdown bank.  It is not the simulator itself; instead it is the place
where the simplified reactor model is derived, calibrated, and sanity-checked
before the resulting constants and algorithms are frozen into the production
code.  It has two estimation levels.

\subsection{Estimate 1: homogeneous buckling sanity check}

The first estimate is a coarse analytical model based on one-speed diffusion
theory in a homogeneous bare cylinder:
\[
-D\nabla^2 \phi + \Sigma_a \phi = \frac{1}{k}\nu\Sigma_f \phi.
\]

The criticality condition is written as
\[
B_m^2 = \frac{\nu\Sigma_f - \Sigma_a}{D} = B_g^2,
\]
with geometric buckling for a finite cylinder approximated by
\[
B_g^2 \approx \left(\frac{2.405}{R}\right)^2 + \left(\frac{\pi}{H}\right)^2.
\]

In the notebook this estimate is used only as a \textbf{first-pass size check}.
It is intentionally approximate because it ignores:
\begin{itemize}
  \item the inner moderator hole,
  \item reflector heterogeneity,
  \item explicit axial reflector slabs,
  \item and the geometric control rod.
\end{itemize}

\subsection{Estimate 2: full 2D $r$--$z$ finite-difference diffusion model}

The second estimate is the notebook's main result.  It solves a one-group
diffusion eigenvalue problem on a cylindrical $r$--$z$ mesh:
\[
-\nabla \cdot \left(D \nabla \phi\right) + \Sigma_a \phi
= \frac{1}{k_{\mathrm{eff}}}\nu\Sigma_f \phi.
\]

This estimate captures:
\begin{itemize}
  \item radial and axial leakage,
  \item the annular fuel geometry,
  \item inner and outer D$_2$O regions,
  \item axial reflector slabs,
  \item and the real rod tip position as insertion changes.
\end{itemize}

The notebook uses this model to:
\begin{enumerate}
  \item search for critical annular geometry,
  \item compute the clean-core eigenvalue,
  \item scan rod insertion versus reactivity,
  \item and convert that scan into the calibration constants later frozen in the
        backend.
\end{enumerate}

\section{Geometry, materials, and assumptions}

\subsection{Core layout}

The modeled reactor is an annular cylindrical core with a central D$_2$O region,
a homogenized fuel annulus, and an outer D$_2$O reflector.  The current frozen
geometry in \texttt{backend/reactor\_backend/}\linebreak\texttt{calibration.py}
is:

\begin{center}
\begin{tabular}{lc}
\toprule
Quantity & Value \\
\midrule
Inner moderator radius $R_{\mathrm{inner}}$ & $80.0$ cm \\
Fuel outer radius $R_{\mathrm{fuel}}$ & $345.6$ cm \\
Reflector outer radius $R_{\mathrm{refl}}$ & $405.6$ cm \\
Active height $H$ & $691.2$ cm \\
Axial reflector thickness $H_{\mathrm{refl}}$ & $60.0$ cm per side \\
Equivalent rod radius $r_{\mathrm{rod}}$ & $50.0$ cm \\
\bottomrule
\end{tabular}
\end{center}

The model is axisymmetric, so the mesh covers only $r \ge 0$ while still
representing the full 3D core through cylindrical volume weighting
$2\pi r\,dr\,dz$.

\subsection{One-group effective constants}

The notebook starts from microscopic cross sections at the 2200 m/s thermal
reference point and then builds \textbf{engineering-effective} one-group
constants.  The production code uses the following frozen values:

\begin{center}
\begin{tabular}{lccc}
\toprule
Region & $D$ [cm] & $\Sigma_a$ [cm$^{-1}$] & $\nu\Sigma_f$ [cm$^{-1}$] \\
\midrule
Inner D$_2$O moderator & $1.25$ & $3.5\times 10^{-4}$ & $0$ \\
Fuel annulus & $0.95$ & $8.5\times 10^{-3}$ & $8.59\times 10^{-3}$ \\
Outer D$_2$O reflector & $1.15$ & $4.0\times 10^{-4}$ & $0$ \\
\bottomrule
\end{tabular}
\end{center}

These are \textbf{homogenized one-group constants}, not a multigroup transport
library.  The fuel values intentionally fold together:
\begin{itemize}
  \item natural uranium fuel,
  \item D$_2$O moderation inside the lattice,
  \item resonance/self-shielding effects,
  \item and spectrum averaging into one effective thermal group.
\end{itemize}

The control bank is represented as an added absorber in the inner moderator
region with
\[
\Delta \Sigma_{a,\mathrm{rod,max}} = 0.25\ \mathrm{cm}^{-1}.
\]
This is likewise an effective one-group quantity rather than a detailed rod-cell
transport model.

\subsection{How the physics was manipulated into these constants}

This project deliberately compresses more detailed reactor physics into a
one-group annular model that can run interactively.  The main reductions are:
\begin{enumerate}
  \item \textbf{Energy collapse.}  Multigroup slowing-down and spectral effects
        are represented by a single thermal group with effective
        $(D,\Sigma_a,\nu\Sigma_f)$.
  \item \textbf{Lattice homogenization.}  The fuel annulus is not modeled as
        explicit fuel pins in moderator.  Instead natural uranium and D$_2$O are
        volume-averaged into one fissile ring region.
  \item \textbf{Leakage folded into effective losses.}  In a one-group model,
        any loss of neutrons to higher-energy groups or to unresolved geometric
        detail is represented as part of an effective absorption-like loss.
  \item \textbf{Equivalent rod model.}  The real control/shutdown system is
        collapsed to one central cylindrical absorber with an effective radius
        and an effective extra absorption $\Delta\Sigma_a$.
  \item \textbf{Calibration by 2D diffusion, not by transport.}  The final rod
        worth curve and critical insertion are taken from repeated diffusion
        eigenvalue solves, then frozen for the point model.
\end{enumerate}

The notebook is explicit about these manipulations.  For example, the D$_2$O
``effective'' $\Sigma_a$ is much larger than the true 2200 m/s absorption of
pure heavy water, because the one-group model uses $\Sigma_a$ as a \emph{net
loss coefficient} after spectral and geometric simplification.  Likewise, the
fuel-annulus $\Sigma_a$ is intentionally inflated relative to a naive
2200 m/s natural-uranium homogenization so the one-group annular model better
matches the critical-size and rod-worth behavior seen in the notebook's 2D
studies.

\subsection{Boundary assumptions}

The diffusion solves use extrapolated zero-flux boundaries rather than placing
the computational boundary exactly at the physical reflector edge.  In code this
is represented by an extrapolation factor of $2.13 D_{\mathrm{refl}}$, so the
radial and axial computational extents are larger than the physical core.

\section{How the notebook feeds the web app}

The notebook does \emph{not} run the web app directly.  Instead it establishes
the physical basis that is later frozen into backend constants and reusable
solver code:

\begin{enumerate}
  \item \textbf{Estimate 1} gives an analytical reference for rough critical size.
  \item \textbf{Estimate 2} performs the annular $r$--$z$ diffusion study that
        supplies the production calibration.
  \item The resulting geometry, material constants, and rod-worth data are copied
        into \texttt{calibration.py}.
  \item The backend then exposes two kinds of runtime models:
        \begin{itemize}
          \item a calibrated point-kinetics model for fast dashboard transients,
          \item and the validated multigroup diffusion service used by the
                frontend Core page.
        \end{itemize}
\end{enumerate}

So the notebook is best viewed as the \textbf{physics derivation and calibration
workspace}, while the Python backend is the \textbf{production implementation}.

\subsection{How the notebook output becomes backend input}

The handoff from notebook to code is not symbolic; it is numerical.  The
production backend lifts the notebook's final values into
\texttt{calibration.py}:
\begin{itemize}
  \item annular geometry dimensions,
  \item one-group material constants,
  \item clean-core $k_{\mathrm{eff}}$ and excess reactivity,
  \item the 11-point rod-worth table,
  \item and the derived critical insertion percentage.
\end{itemize}

This means the backend no longer recomputes those design-calibration studies on
startup.  Instead it assumes the notebook has already answered the
``what reactor are we modeling?'' question and the Python solvers answer the
``how does that reactor evolve for this input?'' question.

\section{Overview page: point-kinetics transient model}

\subsection{Governing equations}

The Overview page uses the standard six-group delayed-neutron point-kinetics
system:
\begin{align}
\frac{dn}{dt} &= \frac{\rho - \beta}{\Lambda} n + \sum_{i=1}^{6}\lambda_i C_i,
\\
\frac{dC_i}{dt} &= \frac{\beta_i}{\Lambda} n - \lambda_i C_i,
\qquad i=1,\dots,6.
\end{align}

Here:
\begin{itemize}
  \item $n(t)$ is the normalized neutron population,
  \item $\rho$ is reactivity in $\Delta k/k$,
  \item $\beta = \sum_i \beta_i$ is the effective delayed-neutron fraction,
  \item $\Lambda$ is the prompt neutron generation time,
  \item $C_i$ is precursor concentration for delayed group $i$,
  \item $\lambda_i$ is the delayed-group decay constant.
\end{itemize}

\subsection{Constants used by the web app}

The current production constants in \texttt{config.py} are loaded from the
published four-group OpenMC MGXS export
\texttt{openmc/reference\_data/concentric/group\_sweep/group\_4/outputs/mgxs\_constants.json}.
The backend averages the fuel-ring delayed-neutron data with a fission-source
weight,
\[
w_j = \sum_g \left(\nu\Sigma_f\right)_{j,g}\Phi_{j,g},
\]
where $j$ indexes the resolved fuel-ring domains.  The fuel-ring $\beta_i$
values are \emph{averaged}, not summed across rings.  The prompt neutron
generation time is taken from the export's \texttt{prompt\_generation\_time\_s}
tally.

The resulting constants are:
\begin{align*}
\beta_{\mathrm{eff}} &= 0.00678910882978 = 678.91\ \mathrm{pcm}, \\
\Lambda &= 0.00417212833712\ \mathrm{s}.
\end{align*}

The six delayed groups are:

\begin{center}
\begin{tabular}{ccc}
\toprule
Group & $\beta_i$ & $\lambda_i$ [s$^{-1}$] \\
\midrule
1 & 0.000227917094163 & 0.0133443203744 \\
2 & 0.00119534670498 & 0.0326776283533 \\
3 & 0.00115198689522 & 0.120914791504 \\
4 & 0.00262514724228 & 0.304202879275 \\
5 & 0.00112075519129 & 0.855415421400 \\
6 & 0.000467955701846 & 2.87295916880 \\
\bottomrule
\end{tabular}
\end{center}

Displayed output is scaled from the nominal operating point
\[
P_0 = 20\ \mathrm{MW_{th}}, \qquad
\Phi_0 = 1.5\times 10^{12}\ \mathrm{n\,cm^{-2}\,s^{-1}},
\]
using
\[
P(t) = P_0 n(t), \qquad \Phi(t) = \Phi_0 n(t).
\]

\subsection{How rod motion enters the point model}

The point-kinetics solver does not compute rod worth from first principles during
runtime.  Instead it loads the OpenMC continuous-energy rod-worth table
\texttt{openmc/reference\_data/concentric/rod\_scan/results/rod\_worth.csv}.
Reactivity is computed as
\[
\rho_{\mathrm{tot,pcm}}(x)
= \rho_{\mathrm{clean,pcm}} + \Delta \rho_{\mathrm{rod,pcm}}(x)
+ \rho_{\mathrm{scram,pcm}},
\]
where $x \in [0,1]$ is rod insertion fraction.  In the current calibration:
\begin{align*}
k_{\mathrm{clean}} &= 1.00224587, \\
\rho_{\mathrm{clean}} &= +224.084\ \mathrm{pcm}, \\
x_{\mathrm{crit}} &\approx 0.265, \\
\Delta \rho_{\mathrm{rod}}(1.0) &= -1350.73\ \mathrm{pcm}, \\
\rho_{\mathrm{scram}} &= -450\ \mathrm{pcm}\ \text{when latched.}
\end{align*}

The backend interpolates linearly between the insertion points stored in the
CSV.  The older diffusion-derived constants remain only as parameters for the
legacy one-group spatial display model.

\subsection{Thermal reactivity feedback}

The point model also adds thermal feedback when the Modelica FMU reports all
three effective feedback signals:
\texttt{T\_fuelEff}, \texttt{T\_moderatorEff}, and \texttt{rho\_m\_eff}.  The
coefficients are loaded from
\texttt{openmc/reference\_data/concentric/reactivity\_coefficients/results/reactivity\_coefficients.csv}:
\begin{align*}
\alpha_f &= -2.138272335\ \mathrm{pcm/K}, \\
\alpha_{m,T} &= -10.72767894\ \mathrm{pcm/K}, \\
\alpha_{m,\rho} &= 20993.04779\ \mathrm{pcm/(g/cm^3)}.
\end{align*}

At reset, the current FMU effective state is stored as the zero-feedback
reference.  Runtime feedback is then
\[
\rho_{\mathrm{fb,pcm}} =
\alpha_f \left(T_{f,eff}-T_{f,0}\right)
+ \alpha_{m,T}\left(T_{m,eff}-T_{m,0}\right)
+ \alpha_{m,\rho}\left(\rho_{m,eff}-\rho_{m,0}\right).
\]
This preserves the existing critical startup behavior while still using the
OpenMC finite-difference slopes.  Thermal feedback is \textbf{FMI-only}: if the
FMU is unavailable and the fallback thermal model is active, the feedback terms
are reported as zero and are not added to reactivity.

\subsection{Numerical method}

The point-kinetics equations are stiff, so the backend does not use explicit
Euler.  Instead \texttt{engine.py} advances both the neutron population and the
precursors implicitly at the new time level and then algebraically eliminates
the precursor unknowns.

Start from the backward-Euler discretization
\begin{align}
\frac{n^{k+1} - n^k}{\Delta t}
&=
\frac{\rho-\beta}{\Lambda} n^{k+1}
+ \sum_{i=1}^6 \lambda_i C_i^{k+1},
\\
\frac{C_i^{k+1} - C_i^k}{\Delta t}
&=
\frac{\beta_i}{\Lambda} n^{k+1} - \lambda_i C_i^{k+1}.
\end{align}

Solving the precursor equation for the new precursor gives
\[
C_i^{k+1}
=
\frac{C_i^k + \Delta t\,(\beta_i/\Lambda)\,n^{k+1}}
{1 + \lambda_i \Delta t},
\]
which is the relation already visible in the source.

Substituting that expression into the neutron balance yields
\[
\frac{n^{k+1} - n^k}{\Delta t}
=
\frac{\rho-\beta}{\Lambda} n^{k+1}
+ \sum_i \lambda_i
\frac{C_i^k + \Delta t\,(\beta_i/\Lambda)\,n^{k+1}}
{1+\lambda_i\Delta t}.
\]

Now separate the terms involving only known old-state data from the terms that
still multiply $n^{k+1}$:
\[
\frac{n^{k+1} - n^k}{\Delta t}
=
\sum_i \frac{\lambda_i C_i^k}{1+\lambda_i\Delta t}
+
\left[
\frac{\rho-\beta}{\Lambda}
+ \sum_i \frac{\lambda_i \Delta t\,\beta_i/\Lambda}{1+\lambda_i\Delta t}
\right] n^{k+1}.
\]

Rearranging gives one scalar solve for the new neutron population,
\[
n^{k+1}
=
\frac{n^k + \Delta t\sum_i \dfrac{\lambda_i C_i^k}{1+\lambda_i\Delta t}}
{1 - \Delta t\,\dfrac{\rho-\beta}{\Lambda}
 - \Delta t\sum_i
   \dfrac{\lambda_i \Delta t\,\beta_i/\Lambda}{1+\lambda_i\Delta t}}.
\]

In the implementation this matches the code variables
\texttt{delayed\_source}, \texttt{delayed\_coupling}, \texttt{numerator}, and
\texttt{denominator} in \texttt{engine.py}.  The precursor update then uses the
newly computed $n^{k+1}$.

Numerically, this matters because the prompt part of the kinetics is very fast
relative to the slowest delayed groups.  A naive explicit step would require a
much smaller $\Delta t$ for stability, while the semi-implicit update remains
well behaved for the chosen browser-friendly step size.

\subsection{How the Overview transient is operationalized}

The point-kinetics equations are not advanced continuously in a background
thread.  Instead \texttt{service.py} advances them \emph{opportunistically}
whenever the frontend polls or sends a command:
\begin{enumerate}
  \item compute elapsed wall time since the last request,
  \item cap that wall-time gap to avoid large jumps after browser pauses,
  \item multiply by the configured time-scale factor,
  \item then march the engine forward in repeated $0.02$ s internal steps.
\end{enumerate}

This architecture is a numerical simplification as much as a software one: the
browser sees a continuous transient, but the backend actually performs a sequence
of bounded substeps and samples history only every $0.25$ s.

The current runtime tuning is:
\begin{center}
\begin{tabular}{lc}
\toprule
Parameter & Value \\
\midrule
Integrator step & $0.02$ s \\
History sampling & $0.25$ s \\
Time scale & $8\times$ wall clock \\
Frontend poll interval & $100$ ms \\
Automatic SCRAM threshold & $26$ MW$_{\mathrm{th}}$ \\
\bottomrule
\end{tabular}
\end{center}

\section{Legacy one-group core backend: steady 2D diffusion solution}

This legacy backend service reuses the same annular one-group model as the
notebook and solves the 2D eigenvalue problem for a supplied rod position.  It
is retained as reference code and for historical calibration context, but it is
no longer exposed as the frontend Core page.

\subsection{Discretization}

The implementation in \texttt{diffusion.py} and \texttt{core\_service.py} uses:
\begin{itemize}
  \item a cell-centered finite-difference mesh in $r$ and $z$,
  \item harmonic averaging of diffusion coefficients at material interfaces,
  \item sparse matrix assembly for leakage, absorption, and fission terms,
  \item and a power-iteration style eigenvalue solve for $k_{\mathrm{eff}}$ and
        the fundamental flux shape.
\end{itemize}

The legacy core solve runs at
\[
\Delta r = \Delta z = 3\ \mathrm{cm},
\]
and results are cached by rod insertion state so repeated requests do not repeat
the full eigenvalue solve.

\subsection{How the steady diffusion matrices are built}

The Python solver does not discretize the Laplacian in Cartesian form and then
``pretend'' the reactor is cylindrical.  It assembles the operator directly in
cylindrical coordinates.  On each cell-centered mesh point $(r_i,z_j)$ it builds:
\begin{itemize}
  \item radial couplings weighted by the cylindrical face areas $r_{i\pm 1/2}$,
  \item axial couplings weighted by $1/\Delta z^2$,
  \item and a diagonal absorption term from the local $\Sigma_a$.
\end{itemize}

Material interfaces are handled with the harmonic mean of adjacent diffusion
coefficients, which is the standard finite-volume-like choice for discontinuous
media because it preserves current continuity better than a simple arithmetic
average.

The assembled linear system can be written schematically as
\[
L\phi = \frac{1}{k}F\phi,
\]
where:
\begin{itemize}
  \item $L$ is a sparse loss matrix containing leakage and absorption,
  \item $F$ is a diagonal sparse matrix containing $\nu\Sigma_f$ in fuel cells
        and zero elsewhere.
\end{itemize}

The reason for that structure is direct.  After cell-centered finite-difference
discretization, the loss term at cell $(i,j)$ depends only on:
\begin{itemize}
  \item the flux in the same cell,
  \item the two radial neighbors $(i-1,j)$ and $(i+1,j)$,
  \item and the two axial neighbors $(i,j-1)$ and $(i,j+1)$.
\end{itemize}

So each row of $L$ contains only a small five-point stencil worth of nonzero
entries, making $L$ sparse.  By contrast, the fission source has no spatial
derivative in this one-group model; it is simply
\[
\left(F\phi\right)_{ij} = \nu\Sigma_{f,ij}\phi_{ij}.
\]
That is a cell-local multiplication, so after flattening the 2D mesh into a
single vector, $F$ becomes diagonal.

In other words:
\begin{itemize}
  \item \textbf{leakage/absorption} move neutrons between nearby cells or remove
        them locally $\rightarrow$ sparse coupling matrix $L$,
  \item \textbf{fission production} multiplies the flux only in the current cell
        $\rightarrow$ diagonal production matrix $F$.
\end{itemize}

\subsection{How the steady eigenvalue solve works numerically}

The function \texttt{solve\_2d()} uses a source-iteration / power-iteration
structure:
\begin{enumerate}
  \item start from a positive flux guess $\phi$ and an initial $k=1$,
  \item compute the fission source $F\phi$,
  \item solve the fixed loss system
        \[
        L\phi^{*} = F\phi
        \]
        with a sparse direct solver,
  \item clip tiny negative roundoff values to zero,
  \item normalize $\phi^{*}$ by its peak,
  \item update the eigenvalue with the ratio of new to old total fission source,
        \[
        k_{\mathrm{new}} = \frac{\sum(F\phi^{*})}{\sum(F\phi)},
        \]
  \item and repeat until $|k_{\mathrm{new}}-k| < \texttt{tol}$ after a few warmup
        iterations.
\end{enumerate}

Two numerical choices are especially important:
\begin{itemize}
  \item the loss matrix $L$ is factorized once per solve with
        \texttt{scipy.sparse.linalg.factorized()}, so each iteration is just a
        triangular solve instead of a fresh matrix inversion;
  \item the flux is normalized by peak value, not by integral power, because the
        steady eigenproblem only needs the shape and eigenvalue.
\end{itemize}

\subsection{How geometry and rod-worth calibration are solved}

The notebook and backend use the steady diffusion solver in two further numerical
loops:
\begin{enumerate}
  \item \textbf{critical-radius search:} bisection on $R_{\mathrm{fuel}}$ until
        the solved $k_{\mathrm{eff}}$ brackets unity to within tolerance;
  \item \textbf{rod-worth scan:} repeated solves at insertion fractions
        $x=0,0.1,\dots,1.0$, converting each $k$ to reactivity
        $\rho = (k-1)/k$ in pcm.
\end{enumerate}

The rod-worth scan may warm-start each solve with the previous flux shape.  That
is a purely numerical acceleration: neighboring rod positions have similar flux
shapes, so reusing the previous converged state reduces iteration count without
changing the underlying model.

\subsection{What the service returns}

The backend returns:
\begin{itemize}
  \item the peak-normalized 2D flux field $\phi(r,z)$,
  \item a radial profile through the midplane,
  \item an axial profile through the fuel-annulus midpoint,
  \item and metadata including $k_{\mathrm{eff}}$, mesh spacing, and iteration count.
\end{itemize}

This was the web app's direct counterpart to the notebook's Estimate 2 flux
maps before the frontend Core page was moved to the validated multigroup
pipeline.

\section{Legacy transient backend: direct time-dependent diffusion}

\subsection{Why this model exists}

The Overview page assumes a fixed spatial shape and evolves only an amplitude.
The legacy transient backend removes that simplification and evolves the
\textbf{spatial flux field itself}, together with the six delayed-neutron
precursor fields.  It is no longer exposed as a frontend page.

\subsection{Governing equations}

The transient diffusion solver in
\texttt{backend/reactor\_backend/transient\_diffusion.py} solves
\begin{align}
\frac{1}{v}\frac{\partial \phi}{\partial t}
&=
-L\phi
+ (1-\beta)\frac{F}{k_{\mathrm{ref}}}\phi
+ \sum_{i=1}^{6}\lambda_i C_i,
\\
\frac{\partial C_i}{\partial t}
&=
\beta_i \frac{\nu\Sigma_f}{k_{\mathrm{ref}}}\phi
- \lambda_i C_i,
\qquad i=1,\dots,6,
\end{align}
where:
\begin{itemize}
  \item $L$ is the spatial loss operator (leakage plus absorption),
  \item $F$ is the diagonal fission-production operator,
  \item $C_i(r,z,t)$ is the precursor field for group $i$,
  \item $v = 2.2\times 10^5\ \mathrm{cm/s}$ is the one-group thermal neutron speed,
  \item $k_{\mathrm{ref}}$ is the eigenvalue of the initialization state.
\end{itemize}

The division by $k_{\mathrm{ref}}$ is deliberate: it references the transient
fission source to the eigenvalue of the initialized state on the transient mesh.
That keeps the transient solve anchored to the same discrete criticality
reference and substantially reduces startup mismatch relative to using the raw
unscaled fission operator.

\subsection{Numerical method}

The transient solver uses \textbf{fully implicit Euler}.  Eliminating the
precursors gives one sparse linear solve per time step:
\[
A\phi^{n+1} = b,
\]
with
\[
A = \frac{1}{v\Delta t}I + L
- b_{\mathrm{eff}}\,\mathrm{diag}\!\left(\frac{\nu\Sigma_f}{k_{\mathrm{ref}}}\right),
\]
and
\[
b_{\mathrm{eff}}
=
(1-\beta) + \sum_i \frac{\beta_i \lambda_i \Delta t}{1+\lambda_i \Delta t}.
\]

After $\phi^{n+1}$ is solved, each precursor field is updated by
\[
C_i^{n+1}
=
\frac{C_i^n + \Delta t\,\beta_i(\nu\Sigma_f/k_{\mathrm{ref}})\phi^{n+1}}
{1+\lambda_i \Delta t}.
\]

This scheme is unconditionally stable for the linear model and is much better
suited than explicit Euler to delayed-neutron reactor kinetics.

In implementation terms, the transient solver performs three numerically distinct
operations:
\begin{enumerate}
  \item \textbf{Initialization:} run a steady \texttt{solve\_2d()} eigenvalue
        solve on the transient mesh, flatten the 2D flux field, and initialize
        each precursor field with the steady relation
        \[
        C_{i,0} = \frac{\beta_i}{\lambda_i}\frac{\nu\Sigma_f}{k_{\mathrm{ref}}}\phi_0.
        \]
  \item \textbf{Matrix rebuild on rod motion:} if rod insertion changes, rebuild
        the sparse matrix $A$ and refactorize it.  In code this cache is keyed by
        the rounded rod fraction together with $k_{\mathrm{ref}}$.
  \item \textbf{Cheap repeated stepping between rod changes:} if rod insertion is
        unchanged, reuse the cached LU factorization so each new time step is
        just one sparse linear solve plus explicit vector updates for the six
        precursor groups.
\end{enumerate}

This split kept the former transient page interactive even though it used a
full spatial solver: the expensive part is the factorization, not the
individual post-factorization time steps.

\subsection{How the transient physics is manipulated into a solvable model}

The legacy transient model is more detailed than point kinetics, but it is still not a
full reactor multiphysics model.  Several physics reductions are built into the
solver:
\begin{itemize}
  \item prompt and delayed neutrons are treated in one energy group,
  \item thermal-hydraulic feedback is absent, so all cross sections stay fixed
        while the transient runs,
  \item rod motion changes only the absorber map inside the central control-bank
        cylinder,
  \item the internally stored flux is an \emph{absolute amplitude field}, while
        UI heatmaps are separately peak-normalized display copies,
  \item and power is derived from a cylindrical fission integral rather than from
        a standalone thermal-hydraulic power balance.
\end{itemize}

That last point is numerically important: the solver never renormalizes its
internal state after each step.  If it did, the transient could not display real
growth or decay in $P/P_0$.

\subsection{Runtime settings}

To keep the direct transient responsive in a browser-driven local app, the
production solver uses:

\begin{center}
\begin{tabular}{lc}
\toprule
Parameter & Value \\
\midrule
Transient mesh & $\Delta r = \Delta z = 5$ cm \\
Time step & $\Delta t = 1$ s \\
Display heatmap & $40 \times 60$ points \\
History limit & 200 points \\
\bottomrule
\end{tabular}
\end{center}

The transient mesh is intentionally coarse so the LU factorization can be
rebuilt quickly whenever rod insertion changes.

\subsection{Service-level time stepping}

The legacy transient backend is advanced by \texttt{transient\_service.py}, which keeps one
shared solver instance and one shared state vector.  On every
\texttt{GET /api/transient-diffusion/state} request it:
\begin{enumerate}
  \item measures elapsed wall time,
  \item converts that to simulated time at a fixed scale of $1\times$,
  \item computes how many whole 1 s transient steps fit in that interval,
  \item caps the advance to at most four steps per request,
  \item and appends one history sample per accepted time step.
\end{enumerate}

This cap is not part of the physics model; it is a numerical and UX safeguard so
one delayed frontend poll cannot trigger an unexpectedly long blocking solve.

\subsection{Power normalization}

The transient diffusion page no longer assumes that power is proportional to one
global amplitude alone.  Instead it evaluates the cylindrical fission-rate
integral
\[
P \propto \sum_j \nu\Sigma_{f,j}\,\phi_j\,V_j,
\qquad
V_j \propto r_j \Delta r \Delta z,
\]
and reports
\[
\frac{P}{P_0}
=
\frac{\sum_j \nu\Sigma_{f,j}\phi_j V_j}
{\sum_j \nu\Sigma_{f,j}\phi_{j,0} V_j}.
\]

The factor $2\pi$ cancels in the ratio, so the implementation only needs the
$r_j \Delta r \Delta z$ weight per cell.

\section{OpenMC-informed resolved multigroup diffusion modeling}

The one-group annular model described above is intentionally simplified for
interactive use.  In parallel, the repository now includes a
\textbf{resolved multigroup diffusion} modeling chain that derives its nuclear
data directly from an OpenMC continuous-energy Monte Carlo transport
calculation.  This section describes that higher-fidelity pipeline.

\subsection{Overview of the pipeline}

The multigroup workflow has four stages, all implemented as reusable Python
modules rather than notebook-only logic:

\begin{enumerate}
  \item \textbf{MGXS export} (\texttt{openmc/mgxs\_export.py}):
        attach multigroup cross-section tallies to each resolved OpenMC cell
        (fuel rings, central channel, coolant, moderator, reflector, control
        rod), run an eigenvalue calculation, and write group constants plus
        delayed-neutron data as JSON and CSV.
  \item \textbf{Adapter and region mapping}
        (\texttt{backend/reactor\_backend/openmc\_mgxs\_adapter.py}):
        parse the OpenMC \texttt{model.xml} to extract cell boundaries
        (z-cylinders and z-planes), load the canonical
        \texttt{mgxs\_constants.json}, validate the scattering contract,
        preserve the centimeter-based OpenMC units, and map each OpenMC cell to a
        \texttt{MultiGroupRegion} inside a \texttt{CylindricalLayeredModel2D}.
  \item \textbf{Finite-volume multigroup diffusion solver}
        (\texttt{backend/reactor\_backend/multigroup\_diffusion.py}):
        assemble per-group sparse spatial operators, solve the multigroup
        eigenvalue problem, and optionally scan rod insertion.
  \item \textbf{CE-referenced SPH qualification}
        (\texttt{backend/reactor\_backend/multigroup\_sph.py}):
        normalize the OpenMC and diffusion region-integrated fluxes to equal
        fission production, fit one factor per resolved region and energy
        group, apply the factors consistently to transport and reaction terms,
        and evaluate eigenvalue, flux, and power-shape acceptance criteria.
        This workflow is currently paused.
  \item \textbf{Persistent preparation cache}
        (\texttt{backend/reactor\_backend/multigroup\_diffusion\_cache.py}):
        fingerprint the MGXS data, mesh settings, SPH artifact, and solver
        configuration; build corrected CSR matrices and a clean-core solution
        once, then persist everything to disk for fast reload.
\end{enumerate}

The command-line entry point is
\texttt{backend/openmc\_to\_diffusion.py}.  The notebook
\texttt{theory/concentricModel.ipynb} remains an analysis interface, while the
web application exposes a four-group rod-aware page backed by the same adapter,
clean-state persistent solution, and in-process rodded solves.

\subsection{MGXS export from OpenMC}

The export module (\texttt{mgxs\_export.py}) takes a previously built OpenMC
\texttt{model.xml} and attaches multigroup tally objects to every resolved
cell domain.  The domain definitions are configurable; the resolved mode maps
each \texttt{fuel\_ring\_N} cell to its own domain label, plus the central
moderator channel, heavy-water coolant, outer moderator, reflector, and
parked control rod.  A frozen \texttt{MGXSExportConfig} dataclass captures the
particle count, batch schedule, energy-group edges (any CASMO structure from
1 to 70 groups), Legendre scattering order (P0, P1, P3, or P5), and the
delayed-neutron group count.

Two scattering representations are tallied from the same continuous-energy
run:
\begin{itemize}
  \item The \textbf{canonical diffusion export} uses non-$\nu$, uncorrected
        scattering (\texttt{scatter\_correction = null}).  The adapter prefers
        the \texttt{consistent scatter matrix} and falls back to
        \texttt{scatter matrix}.  Only the zeroth Legendre moment is passed to
        diffusion; higher moments are recorded and discarded.
  \item A \textbf{supplemental P0-corrected consistent scattering library}
        is tallied for OpenMC's own multigroup transport validation run.
        This correction is validation-only and is not written into the
        canonical diffusion JSON.
\end{itemize}

When \texttt{validation\_scatter\_correction = P0}, the export contract requires
\texttt{legendre\_order = 0}, because OpenMC does not apply that correction to
a higher-order Legendre library.  A higher-order canonical export can still be
read by the adapter, but the validation correction must be disabled for that
run and the diffusion path will retain only moment zero.

The diffusion coefficient is taken directly from OpenMC's transport
cross-section tally:
\[
D_g = \frac{1}{3\,\Sigma_{\mathrm{tr},g}} .
\]

This distinction is important.  OpenMC's P0 scatter correction modifies the
isotropic scattering data used by its multigroup \emph{transport} calculation
to represent part of the first angular moment.  The diffusion calculation
instead represents anisotropic migration through $\Sigma_{\mathrm{tr},g}$ in
$D_g$.  Its scattering source therefore uses the uncorrected non-$\nu$ P0
reaction-rate matrix.  The two calculations use different angular closures;
agreement of the OpenMC multigroup validation does not imply that the
diffusion approximation has been calibrated.

All neutronics quantities remain in OpenMC's native centimeter-based units:
$D$ in cm, macroscopic cross sections in cm$^{-1}$, and geometry and mesh
dimensions in cm.  No SI conversion is performed in the adapter.  This gives
$D\nabla^2\phi$ and $\Sigma\phi$ consistent units throughout the steady
solver.

Each export run produces:
\begin{itemize}
  \item \texttt{mgxs\_constants.json} --- nested group constants, delayed-neutron
        data, metadata for every domain, raw region/group CE flux tally means
        and standard deviations, and a cylindrical CE
        \texttt{kappa-fission} power-mesh tally;
  \item \texttt{group\_constants.csv} --- flat table of all group-constant rows;
  \item \texttt{delayed\_neutrons.csv} --- flat table of all delayed-neutron rows.
\end{itemize}

The raw CE flux tally is retained in fast-to-thermal group order without the
independent groupwise normalization used by the GeN-Foam auxiliary export.
This distinction is necessary because SPH requires relative region/group
reaction rates from one common Monte Carlo normalization.  The cylindrical
power tally supplies independent normalized radial and axial shape checks.

The notebook \texttt{openmc/exportMGXS.ipynb} runs group-count and
Legendre-order sweeps, storing each result under
\texttt{group\_sweep/group\_<N>/}.  It also runs an OpenMC multigroup
validation step that rebuilds an \texttt{mgxs.h5} library and solves the same
geometry in multigroup transport mode, producing the $\Delta\rho$ (MG $-$ CE)
values reported in the workflow document.

\subsection{Adapter: from OpenMC cells to diffusion regions}

The adapter (\texttt{openmc\_mgxs\_adapter.py}) is the bridge between the
OpenMC export and the diffusion solver.  It performs five operations:

\begin{enumerate}
  \item \textbf{Geometry extraction.}  Parses the OpenMC \texttt{model.xml}
        with Python's \texttt{ElementTree}, identifies every named cell, and
        extracts its radial and axial bounds from the z-cylinder and z-plane
        surfaces referenced in the cell's region specification.
  \item \textbf{Contract validation.}  Requires
        \texttt{scatter\_correction = null} in the export config so the P0
        scattering diagonal is not transport-corrected.  Rejects exports that
        are missing required domains (fuel rings, central channel, moderator,
        reflector) or whose MGXS domain labels do not match the XML cell names.
  \item \textbf{Scattering-matrix selection.}  Prefers
        \texttt{consistent scatter matrix}; falls back to \texttt{scatter matrix}.
        Discards any Legendre moments above P0 and interprets the stored matrix
        as $\Sigma_{s,g'\rightarrow g}$, with the first group index denoting the
        source group and the second the destination group.
  \item \textbf{Diffusion-coefficient validation.}  For every active
        region/group pair, checks that the exported coefficient is finite,
        positive, and consistent with
        $D_g=1/(3\Sigma_{\mathrm{tr},g})$ to numerical precision.
  \item \textbf{Inactive-group sanitization.}  Energy groups with zero
        absorption, zero fission, zero chi, and zero scatter coupling are
        flagged as inactive.  Their undefined diffusion coefficients are
        replaced with a finite coefficient from another active group so the
        spatial operator remains well-posed; every replacement is reported in
        the adapter summary.
\end{enumerate}

The output is a \texttt{CylindricalLayeredModel2D} containing one
\texttt{MultiGroupRegion} per OpenMC cell.  Each region holds per-group
vectors for $D$, $\Sigma_a$, $\nu\Sigma_f$, $\kappa\Sigma_f$, $\chi$, and a
$G\times G$ P0 scattering matrix $\Sigma_{s,g'\to g}$.  The adapter also
loads the optional raw CE region/group flux and power-mesh reference.  The
radial mesh is fitted exactly
to all XML-derived material boundaries so no interface is lost to grid
coarsening.  ``Resolved'' here means that no new cross-cell homogenization is
introduced: each exported cell receives its own constants.  The constants are
still flux-weighted averages within that OpenMC tally cell, so this is not a
pin-resolved or continuously varying material description.

\subsection{Finite-volume multigroup diffusion solver}

The solver in \texttt{multigroup\_diffusion.py} solves the $G$-group
steady-state diffusion eigenvalue problem on a 2D cylindrical $(r,z)$ mesh:

\[
-\nabla \cdot \bigl(D_g \nabla \phi_g\bigr)
+ \Sigma_{a,g}\,\phi_g
+ \sum_{g'\neq g} \Sigma_{s,g\to g'}\,\phi_g
\;=\;
\frac{\chi_g}{k_{\mathrm{eff}}}
\sum_{g'} \nu\Sigma_{f,g'}\,\phi_{g'}
\;+\;
\sum_{g'\neq g} \Sigma_{s,g'\to g}\,\phi_{g'} .
\]

The left-hand side collects leakage, absorption, and outscatter from group
$g$ into a single loss operator $A_g$.  The right-hand side contains the
fission source (scaled by $1/k_{\mathrm{eff}}$) and in-scatter from all other
groups.  Within-group scattering does not appear explicitly because its loss
and gain terms cancel in the group balance.  Thus the removal coefficient used
in $A_g$ is
\[
\Sigma_{r,g} =
\Sigma_{a,g} + \sum_{g'\ne g}\Sigma_{s,g\to g'} .
\]

\subsubsection{Spatial discretization}

The mesh is a nonuniform, cell-centered finite-volume grid in cylindrical
coordinates.  Every material radius and axial plane extracted from
\texttt{model.xml} is inserted as a mesh face before the intervals are
subdivided.  The target radial spacing is region dependent (fuel, core
coolant, moderator, and reflector), while one target axial spacing is
currently used throughout.

Integrating the diffusion equation over cell $(i,j)$ gives a balance between
currents through four faces and volumetric reaction terms.  Omitting the
common factor $2\pi$, the cell volume is
\[
V_{ij}' =
\frac{1}{2}\left(r_{i+1/2}^{2}-r_{i-1/2}^{2}\right)\Delta z_j .
\]
For two adjacent cells $P$ and $N$, the face conductance is
\[
H_{PN}
=
\left(
\frac{\delta_P}{D_P}+\frac{\delta_N}{D_N}
\right)^{-1},
\]
where $\delta_P$ and $\delta_N$ are the center-to-face distances.  This is the
harmonic interface treatment implied by continuity of normal current.  After
division by $V_{ij}'$, the off-diagonal coefficient is the negative face
conductance multiplied by the cylindrical face area.  For example,
\[
C_{i+1/2,j}
=
\frac{r_{i+1/2}H_{i,i+1}}
{\tfrac12(r_{i+1/2}^{2}-r_{i-1/2}^{2})},
\]
and the corresponding axial coefficient is
$C_{i,j+1/2}=H_{j,j+1}/\Delta z_j$.  The diagonal contains the sum of all
outgoing leakage coefficients plus $\Sigma_{r,g}$.

The assembled per-group loss operator $A_g$ is a sparse CSR matrix with a
five-point stencil in $(r,z)$.  It is factorized once per solve with a sparse
direct solver (SuperLU via \texttt{scipy.sparse.linalg.factorized}).

\subsubsection{Production-mesh qualification}

SPH constants are tied to a particular spatial homogenization and mesh, so the
mesh is selected before fitting factors.  The current 4-group study compared
the reference spacing
\[
(\Delta r_{\mathrm{fuel}},\Delta r_{\mathrm{cool}},
 \Delta r_{\mathrm{mod}},\Delta r_{\mathrm{refl}},\Delta z)
=(0.1,1,5,10,10)\ {\rm cm}
\]
with the candidate
\[
(0.5,1,5,10,20)\ {\rm cm}.
\]
Both meshes preserve every material face exactly.  The reference contains
15\,096 cells and solved in 1.75~s in the recorded run; the candidate contains
6\,360 cells and solved in 0.42~s.  Although the eigenvalue changed by only
$+5.52$~pcm, the candidate failed the field criteria: its maximum resolved
region/group flux error was 35.3\% in the parked-control-rod region, and its
radial and axial power-shape RMS errors were 4.30\% and 3.07\%.  The required
limits were 0.5\% and 1\%, respectively.  Therefore the 15\,096-cell mesh is
retained for the first production factor set.

Power profiles are compared as power per unit radial or axial coordinate
rather than raw power per mesh bin.  Otherwise changing the bin width would
appear as a physical shape change.

\subsubsection{Boundary conditions}

At $r=0$, the missing west-face term imposes the cylindrical symmetry
condition
\[
\left.\frac{\partial\phi_g}{\partial r}\right|_{r=0}=0 .
\]
At the outer radial, upper axial, and lower axial boundaries, the solver
imposes zero scalar flux at the computational boundary using the half-cell
distance from the last cell center to the boundary face.  The computational
domain is extended beyond the physical reflector by
\[
d_{\mathrm{ext}} = 2.13\,\max_g D_{\mathrm{refl},g},
\]
radially and at each axial end, and the extension is assigned reflector
properties.  Consequently the implemented condition is an approximate
extrapolated-vacuum boundary:
\[
\phi_g=0
\quad\text{at the extrapolated boundary.}
\]
The same extrapolation distance is used for every group because it is based on
the maximum reflector diffusion coefficient.  This is a pragmatic closure,
not a group-dependent transport boundary condition; its sensitivity should be
separated from spatial-mesh and MGXS-condensation errors.

\subsubsection{Eigenvalue iteration}

The multigroup eigenvalue problem is solved by nested iteration:

\begin{enumerate}
  \item \textbf{Outer power iteration.}  Start from a positive flux guess
        and an initial $k_{\mathrm{eff}} = 1$.  At each outer iteration:
        \begin{enumerate}
          \item compute the fission density $f = \sum_g \nu\Sigma_{f,g}\,\phi_g$;
          \item form the fixed fission source $s^{\mathrm{fiss}}_g
                = \chi_g\,f / k_{\mathrm{eff}}$;
          \item solve the coupled scattering system
                $A_g \phi_g = s^{\mathrm{fiss}}_g
                + \sum_{g'\neq g} \Sigma_{s,g'\to g}\,\phi_{g'}$
                for all groups.
        \end{enumerate}
  \item \textbf{Inner scattering iteration.}  The coupled system first receives
        at most five Gauss--Seidel sweeps over destination groups.  The
        requested inner tolerance is scheduled from $10^{-2}$ in early outer
        iterations to the user-specified \texttt{inner\_tol}.
  \item \textbf{Matrix-free GMRES fallback.}  If those Gauss--Seidel sweeps do
        not satisfy the active inner tolerance, the code solves the coupled
        scattering system with matrix-free GMRES, preconditioned by the same
        per-group SuperLU factors.  This avoids assembling a monolithic
        $(G \times N_{\mathrm{cells}}) \times (G \times N_{\mathrm{cells}})$
        matrix, which would produce prohibitive fill-in for higher group
        counts and fine meshes.
  \item \textbf{Eigenvalue update.}  After each outer iteration the eigenvalue
        is updated from the ratio of new to old total fission production,
        \[
        k_{\mathrm{new}}
        = k_{\mathrm{old}}\,
        \frac{\sum_{i,j} V_{ij}\,\sum_g \nu\Sigma_{f,g,ij}\,\phi_{g,ij}^{\mathrm{new}}}
             {\sum_{i,j} V_{ij}\,\sum_g \nu\Sigma_{f,g,ij}\,\phi_{g,ij}^{\mathrm{old}}} .
        \]
  \item \textbf{Flux normalization.}  The flux is normalized so the
        volume-integrated fission production rate equals unity,
        \[
        \sum_{i,j} V_{ij}\,\sum_g \nu\Sigma_{f,g,ij}\,\phi_{g,ij} = 1 .
        \]
        This fixes the otherwise arbitrary eigenvector scale and improves
        numerical conditioning.  It is not an absolute power normalization;
        conversion to physical flux or power requires a separate amplitude
        based on reactor power and recoverable energy per fission.
\end{enumerate}

Convergence is declared only after the final requested inner tolerance is
active and both
\[
\frac{|k_{\mathrm{new}}-k_{\mathrm{old}}|}
{|k_{\mathrm{new}}|}<\texttt{tol}
\]
and the volume-weighted change in fission-source shape is below
\texttt{source\_tol}.  A volume-weighted equation-balance residual is computed
after convergence and reported as a diagnostic, but it is not currently a
stopping criterion.

\subsection{CE-referenced superhomogenization}

The SPH implementation uses one positive factor $\mu_{i,g}$ for each resolved
OpenMC region $i$ and energy group $g$.  Raw CE and diffusion
region-integrated scalar fluxes are first normalized so that both have unit
total fission production:
\[
\sum_{i,g}\nu\Sigma_{f,i,g}\Phi_{i,g}=1.
\]
For the current iterate, the direct fixed-point target is
\[
\mu_{i,g}^{*}
=
\frac{\Phi_{i,g}^{\mathrm{CE}}}
     {\Phi_{i,g}^{\mathrm{diff}}}.
\]
The code damps this target in logarithmic space,
\[
\ln\mu_{i,g}^{(n+1)}
=(1-\alpha)\ln\mu_{i,g}^{(n)}
+ \alpha\ln\mu_{i,g}^{*},
\qquad \alpha=0.5 ,
\]
which preserves positivity and avoids an asymmetric additive update.

For each source group, the factor multiplies absorption, fission,
\texttt{kappa-fission}, and the complete scattering row:
\[
\widetilde{\Sigma}_{a,i,g}=\mu_{i,g}\Sigma_{a,i,g},\qquad
\widetilde{\nu\Sigma}_{f,i,g}=\mu_{i,g}\nu\Sigma_{f,i,g},
\]
\[
\widetilde{\Sigma}_{s,i,g\rightarrow g'}
=\mu_{i,g}\Sigma_{s,i,g\rightarrow g'} .
\]
The transport term is transformed by the same factor.  Since the solver stores
$D=1/(3\Sigma_{\mathrm{tr}})$, this becomes
\[
\widetilde{D}_{i,g}=\frac{D_{i,g}}{\mu_{i,g}} .
\]
The fission spectrum $\chi$ is unchanged.  Thus the fitting objective is flux
and reaction-rate equivalence, not direct adjustment of $k_{\mathrm{eff}}$.

Only statistically active bins enter convergence and qualification.  The fit
requires active transformed flux errors below 0.2\% and an eigenvalue change
below 5~pcm for two successive iterations; bounded factors and an iteration
limit make failure explicit.  The immutable \texttt{SphFactorSet} records the
factor matrix, active mask, convergence history, mesh, source fingerprint,
algorithm version, and provisional state.  A factor artifact is rejected when
the MGXS JSON, XML geometry, mesh, or algorithm version changes.

Final qualification additionally requires a CE eigenvalue uncertainty no
larger than 15~pcm, corrected $k_{\mathrm{eff}}$ within 50~pcm of the CE mean,
all statistically active region/group fluxes within 1\%, and normalized radial
and axial power profiles within 2\% RMS and 5\% maximum pointwise error.
Repeated corrected solves must agree within 1~pcm.

\subsubsection{Post-solve power-shape correction}

The current production correction is not SPH.  It is a post-solve power-map
correction used only after the diffusion flux and eigenvalue have been
computed.  The clean continuous-energy kappa-fission mesh and the clean
diffusion power map are normalized by total power before forming
\[
C_{\mathrm{clean}}(r,z)
= \frac{\widehat{P}_{\mathrm{CE}}(r,z;0)}
       {\widehat{P}_{\mathrm{diff}}(r,z;0)} .
\]
For a later rod position $x$, the corrected map keeps this factor fixed:
\[
\widehat{P}_{\mathrm{map}}(r,z;x)
= \operatorname{normalize}\left[
  C_{\mathrm{clean}}(r,z) P_{\mathrm{diff}}(r,z;x)
\right].
\]
This is equivalent to applying the relative rod-induced shape change predicted
by diffusion to the clean CE power shape.  The heat-source map used by a
point-kinetics coupling is then
\[
q'''(r,z,t)
= P_{\mathrm{PK}}(t)\widehat{P}_{\mathrm{map}}(r,z;x(t)).
\]
The correction does not change MGXS, flux, or $k_{\mathrm{eff}}$.

\subsubsection{Rod-worth scan}

The equivalent-absorber rod model adds a scalar $\Delta\Sigma_a$ inside the
physical rod radius above an axially moving rod tip.  In those cells it also
sets $\nu\Sigma_f$, $\chi$, and scattering to zero.  This is an equivalent
black/absorbing-region perturbation, not a material replacement using the
parked-rod MGXS.  A rod scan sweeps the insertion fraction
$x \in [0,1]$, rebuilding and factorizing the affected group operators at each
position.  The previous converged flux may be used as the next initial guess.

The resulting $\Delta\rho(x)$ is explicitly exploratory.  Backend point
kinetics uses the OpenMC CE rod-worth CSV instead.  The equivalent absorber is
kept because it gives a practical rodded diffusion power shape, but it is not a
material replacement using rodded MGXS.  The present limitations are that the
available MGXS were generated for the parked rod geometry, the physical rod is
shorter than the desired full-height absorber representation, and the absorber
composition used for the OpenMC worth scan was intentionally weakened relative
to default B4C.  Clean-core SPH factors, when available, do not qualify this
rod perturbation.

\subsection{Persistent preparation cache}

The full solve chain (adapter $\to$ mesh construction $\to$ per-group CSR
assembly $\to$ convergence) is expensive, especially for higher group counts
and fine meshes.  The cache module
(\texttt{multigroup\_diffusion\_cache.py}) avoids recomputing this work on
every notebook restart or explicit cache preparation.

The cache fingerprint is a SHA-256 hash of:
\begin{itemize}
  \item the canonical \texttt{mgxs\_constants.json} contents,
  \item the colocated \texttt{model.xml} contents,
  \item the mesh-spacing configuration (target $\Delta r$ and $\Delta z$ per
        region),
  \item the complete SPH factor artifact, including its algorithm version,
        when corrections are applied,
  \item the rod-absorption parameter and solver tolerances,
  \item and an explicit cache-schema version.
\end{itemize}

The fingerprint does not hash the Python source itself.  A numerical algorithm
change that affects cached data therefore requires incrementing the cache
schema version.

On a cache hit, the following are loaded directly from disk:
\begin{itemize}
  \item radial and axial mesh edge arrays,
  \item cell-wise material arrays ($D$, $\Sigma_a$, $\nu\Sigma_f$,
        $\kappa\Sigma_f$, $\chi$, P0 scattering, and resolved-region indices),
  \item one CSR spatial matrix per energy group,
  \item and the converged clean-core eigenvalue, flux field, and power density.
\end{itemize}

SuperLU factorization objects are process-local and are not serialized; a
process that needs to solve again (e.g.\ for a rod scan) loads the CSR
matrices and re-factorizes them in memory.  The persisted operators describe
the clean state only.  A changed rod position alters the absorption and
scattering maps, so its operators must still be rebuilt and factorized.  The
web service keeps those rodded multigroup solves in an in-process cache keyed
by insertion fraction, but they are not written to the persistent preparation
cache.  On a cache miss, the full clean-state build-and-solve pipeline runs
synchronously and the results are persisted.

The default cache root is \texttt{backend/.cache/openmc\_diffusion/}.  The
CLI entry point \texttt{backend/openmc\_to\_diffusion.py} exposes a
\texttt{--qualify-mesh} study, \texttt{--fit-sph} factor generation, and
\texttt{--prepare-cache} operator preparation for use outside the notebook.

\subsection{Integration with the web application}

The web backend now exposes the four-group rod-aware model through
\texttt{GET /api/multigroup-diffusion/state} and
\texttt{POST /api/multigroup-diffusion/recompute}.  The service is initialized
lazily.  Its default path intentionally matches
\texttt{openmc/diffusion\_concentric\_reactor.ipynb}: four energy groups,
the \texttt{groupwise\_fvm} boundary treatment, 0.1/1/5/20 cm radial mesh
targets for fuel/coolant/moderator/reflector, and a 10 cm axial target.  The
clean unrodded solve gives
$k_{\mathrm{eff}}=1.001429$ ($+142.7$ pcm relative to critical), within about
$1$ pcm of the clean continuous-energy OpenMC reference.

The flux map shown by the frontend is the raw four-group diffusion scalar flux,
mirrored from cylindrical $r$-$z$ into an axisymmetric $x$-$z$ image.  Power is
displayed with the fixed clean CE power-shape factor when a CE power mesh is
available.  Each response reads the current rod insertion from the
point-kinetics simulation, solves or reuses the corresponding rodded
multigroup diffusion state, and applies the fixed clean factor to the rodded
diffusion power map.  The frontend names this view \textbf{Core}; the backend
route names remain \texttt{multigroup-diffusion} for compatibility.

SPH fitting and qualification remain available in the offline workflow, but
the online Core page does not apply SPH factors.  Rodded maps are marked
provisional because they use an equivalent absorber, not rodded MGXS or
rod-dependent transport corrections.  Reactivity remains tied to the OpenMC CE
rod-worth CSV, not to the equivalent diffusion rod-worth scan.

\subsection{Current status and verification}

The adapter and export configuration support the listed 1, 2, 4, 8, 16, 25,
40, and 70 group structures.  The documented diffusion benchmark table
currently contains resolved clean-core results through 16 groups; support for a
group structure should not be confused with a completed mesh/runtime
verification at that group count.  The workflow document
\texttt{theory/diffusion\_mgxs\_workflow.md} records the current comparison:

\begin{itemize}
  \item OpenMC multigroup transport with the validation-only P0 correction
        approaches the continuous-energy reference as the group structure is
        refined: the reported MG biases are $+806$, $+426$, $+166$, and
        $+66$~pcm for 2, 4, 8, and 16 groups, respectively.
  \item The uncorrected diffusion bias (diffusion $k_{\mathrm{eff}}$ minus
        OpenMC MG $k_{\mathrm{eff}}$) stabilizes around $-400$ to $-500$~pcm
        for the documented 4, 8, and 16 group cases.  This is evidence of a
        repeatable residual bias, not proof that diffusion closure alone is
        its cause.
  \item The resolved two-group diffusion result is anomalous relative to that
        trend: its bias is about $+1194$~pcm relative to two-group OpenMC MG.
        This points to sensitivity to the very broad energy condensation and
        groupwise diffusion coefficients rather than a simple monotonic
        group-count convergence law.
  \item The resolved one-group diffusion result carries an approximately
        $+1800$~pcm bias relative to continuous energy.  Because the fuel rings
        remain spatially resolved in this calculation, that result cannot be
        attributed to collapsing the rings into one spatial region.
\end{itemize}

The SPH fitting, application, fingerprinting, cache, and qualification
machinery are implemented and unit tested, but the workflow is paused.  The
trial axial flux and power-shape SPH variants are archived for later
investigation because they did not yet produce a qualified correction.  The
online four-group result therefore uses the post-solve CE power-map correction
described above and remains explicitly provisional until a new high-statistics
CE export and a qualified SPH factor set are available.  Discontinuity factors
and true rodded MGXS remain deferred.

\section{How the modeling layers relate}

\begin{center}
\begin{tabular}{p{3.2cm}p{4.0cm}p{5.3cm}}
\toprule
Layer & Model & Purpose \\
\midrule
OpenMC MGXS export & Resolved-cell multigroup & Generate transport-derived group constants and scattering matrices from continuous-energy Monte Carlo \\
\rule{0pt}{13pt}%
Core page / multigroup diffusion solver & $G$-group 2D finite-volume with optional SPH & Solve the resolved-cell flux shape; display rodded power maps with a fixed clean CE power-shape correction \\
Notebook estimate 1 & Homogeneous buckling & First-pass analytical check of scale and leakage trends \\
Notebook estimate 2 / legacy backend core model & Steady 2D one-group diffusion & Compute $k_{\mathrm{eff}}$, flux maps, and historical rod-worth calibration from effective constants \\
Overview page & 0D six-group point kinetics & Fast transient driven by the calibrated reactivity curve \\
Legacy backend transient model & 2D diffusion + 6 precursor fields & Reference spatial transient model, not exposed in the frontend \\
\bottomrule
\end{tabular}
\end{center}

The conceptual workflow is therefore:
\begin{enumerate}
  \item Run OpenMC with resolved-cell MGXS tallies; export constants as JSON,
  \item parse the export through the adapter and build a boundary-fitted
        finite-volume mesh,
  \item qualify the production mesh and archive SPH trials until a qualified
        CE-referenced factor set is available,
  \item solve the multigroup diffusion eigenvalue problem through the CLI,
        notebook, or rod-aware web page,
  \item use the fixed clean CE power-shape factor for rodded display maps
        until rodded MGXS and transport corrections exist,
  \item use the OpenMC rod-worth CSV for point-kinetics reactivity,
  \item use the existing fast point model for dashboard-style operation on the
        Overview page,
  \item and use the diffusion transient (one-group or, in the future,
        multigroup) when spatial dynamics matter.
\end{enumerate}

\section{Important assumptions and limitations}

The one-group web-application models remain intentionally simplified:
\begin{itemize}
  \item one energy group only,
  \item no thermal-hydraulic feedback,
  \item no fuel temperature, moderator temperature, void, xenon, or burnup effects,
  \item homogenized fuel and rod regions,
  \item heavy-water moderation is represented through effective diffusion and
        absorption constants rather than coupled moderator state equations.
\end{itemize}

The OpenMC-informed multigroup diffusion pipeline lifts the one-group
limitation (the configured CASMO structures contain up to 70 groups) and
avoids any additional cross-cell homogenization by preserving every exported
OpenMC cell as a separate diffusion region.  It still uses flux-weighted
cell-homogenized constants and does not yet include:
\begin{itemize}
  \item a qualified high-statistics clean-core SPH factor artifact (the
        fitting and application machinery is implemented),
  \item discontinuity factors,
  \item rodded MGXS or rod-dependent SPH factors,
  \item thermal-hydraulic or burnup feedback,
  \item time-dependent multigroup solves in the web backend.
\end{itemize}

So the repository should be read as an \textbf{educational reactor simulator with
increasing levels of neutronics fidelity}, not as a licensing-grade core code.

\section{Bottom line}

The repository currently implements a hierarchy of reactor models with
increasing spatial and spectral fidelity:

\begin{itemize}
  \item At the top, \textbf{OpenMC continuous-energy Monte Carlo} provides
        reference $k_{\mathrm{eff}}$ values and generates resolved-cell
        multigroup cross sections.
  \item A \textbf{resolved multigroup diffusion solver} consumes those MGXS
        exports, solves the $G$-group eigenvalue problem on a
        boundary-fitted cylindrical mesh, and produces clean or rodded flux
        shapes.  Power maps use a fixed clean CE power-shape correction.  The
        current stored export has not yet met the requirements for a qualified
        SPH factor set.
  \item The older \textbf{one-group annular diffusion model} remains as
        backend reference code only.  It is no longer exposed as a frontend
        Core or Transient Diffusion page.  Point-kinetics rod worth now comes
        from the OpenMC CE rod scan.
  \item The \textbf{point-kinetics transient} on the Overview page gives
        fast control-response exploration in the browser.
  \item The \textbf{direct diffusion transient} remains a backend reference
        model for evolving spatial flux shapes during rod motion.
\end{itemize}

The CLI \texttt{backend/openmc\_to\_diffusion.py}, the notebook
\texttt{theory/concentricModel.ipynb}, and the frontend Core page
are user-facing interfaces to the same multigroup pipeline.  The notebook
retains the equivalent-rod scan for shape exploration and comparison against
the OpenMC rod-worth CSV.  The
\texttt{diffusion\_mgxs\_workflow.md} document records the group-by-group
verification results, mesh study, SPH equations, and scattering contract that
tie the export format to the solver expectations.

That is the current modeling hierarchy of the project.

\end{document}
